The hit rate is low. RBX1 was picked to be hard, and 2.8\% reflects
that. The two prior Adaptyv community campaigns (EGFR 2024 at
\textasciitilde21\%, Nipah-G 2025--26 at \textasciitilde9\% on the
public collection) used larger, more druggable targets with deeper prior
art. RBX1 has neither --- the disordered N-terminal tail offers no
stable interface, and the RING core depends on metal coordination that
the default-buffer screen does not enforce.

The methods that produced binders are not a clean ranking. RFdiffusion
made the most submissions but did not place a binder at the top; ORBIT,
an unpublished stack from Mandrake Bio, placed the Strong hit with a
single 100-aa miniprotein. BindCraft 2 {[}@pacesa2025bindcraft{]}
produced three Medium binders from seven submissions --- a 43\% hit rate
on its own slice, but a slice too thin for a significance claim. AF2
hallucination with ADFlip produced the two next-tightest binders. The
honest summary is that several distinct philosophies of design --- full
diffusion, AF2 backprop with a sequence editor, composite-objective JAX,
and the more recent inverse-folding-style approaches --- each delivered
a top-five binder.

Mini binders dominate. We do not read this as evidence that mini binders
are intrinsically better against RBX1; instead, we read it as evidence
that the field is biased toward mini binders, because most publicly
available design tools target the modality. The submission composition
(25\% miniprotein, 13\% antibody-family, 10\% peptide, 51\%
\emph{Other}) reflects what people built rather than what worked.

What the dataset is useful for: the 322 sequences plus 322 ESMFold
predictions plus 943 SPR replicate traces are a public benchmark for
model calibration on a hard target. The most informative follow-up is a
Zn²⁺-loaded rerun --- the cleanest available probe of whether designs
are targeting the apo or the holo conformation.

We do not claim RBX1 is now a designable target. With one binder under
100 nM, two between 100 and 200 nM, and the rest sitting between 200 nM
and 500 nM, the community has not yet shown that design pipelines
generalise to small, metal-coordinated, partially disordered targets. We
claim, more cautiously, that this 322-design dataset exposes how thin
the binder pool is at this difficulty, and where the next round of
follow-ups should land.
